\scanpage{3}\mspage{2}
% slanted
\margin{b) se déduit de a) et de 1.2. a)}
\diagram{\begin{tikzcd}[ampersand replacement=\&,column sep=small,row sep=small,
execute at end picture={\draw (aa.north west) .. controls +(0,-.25) and +(.04,.16) .. (aa.south west) .. controls +(-.04,-.10) and +(-.16,.03) .. (bb.north west) .. controls +(.08,-.10) and +(-.02,.25) .. (bb.south west);}]
{\text{(1.3)}} \& |[alias=aa]| {\text{a)}} \& {\mathcal{F}\hookrightarrow\unsure{\mathcal{P}|\mathscr{P}}(\mathcal{M})} \& {F\mapsto\widetilde F=\operatorname{\text{Dépl}}(F)\subset\mathcal{M}} \\
\& |[alias=bb]| {\text{b)}} \& {\mathcal{F}\hookrightarrow\mathcal{P}(\mathcal{P}(\mathcal{M}))} \& {\text{plus précisément }\mathcal{F}\to\operatorname{Fig}(\mathcal{M})} \\
\& {\substack{\text{ne sert}\\\text{à rien}}}\arrow[u,bend left=24]\arrow[rr,bend right=13] \& \& {F\mapsto\deleted{\unsure{\sim}}\operatorname{Fig}_{\mathcal{M}}(F)\subset\mathcal{P}(\mathcal{M})}
\end{tikzcd}}
où $\operatorname{Fig}(\mathcal{M})$ désigne la partie de $\mathcal{P}(\mathcal{P}(\mathcal{M}))$ formée des figures ensemblistes dans $\mathcal{M}$. \deleted{\unsure{On}}

On peut considérer que la première application correspond à une relation entre $\mathcal{M}$ et $\mathcal{F}$, appelée relation d'incidence. Pour une figure $F$, $\widetilde F$ s'appelle l'ensemble des \uline{multistrates} \uline{incidentes}, ou le \uline{déployement} de la figure $F$. Si $X\in\mathcal{M}$, $F\in\mathcal{F}$, on dira que la multistrate $X$ est \uline{incidente} à la figure $F$, ou encore que c'est une \uline{strate de} \deleted{$F$}\inserted{\uline{la figure} $F$}, si $X\in\widetilde F$. D'autre part, tout\deleted{\unsure{e}} élément \inserted{$X$} de $\mathcal{M}$ (i.e. \deleted{\unsure{toute}} multistrate), regardé comme une figure par (1.2), admet un déployement $\widetilde X$, et \deleted{la multistrate} \deleted{\unsure{associée}} \deleted{\unsure{on}} on pose
\margin{$\unsure{\bullet|\cdot}$}
\[
\begin{aligned}
\operatorname{Fig}_{\mathcal{M}}(F)&=\{\widetilde X\mid X\in\widetilde F=\operatorname{Depl}(F)\}\subset\mathcal{P}(\mathcal{M})\\
\text{ie.}\quad \operatorname{Fig}_{\mathcal{M}}(F)&\in\mathcal{P}(\mathcal{P}(\mathcal{M}))
\end{aligned}\tag*{(1.4)}
\]
et il résultera des axiomes que c'est une figure ensembliste dans $\mathcal{M}$, \unsure{indexée} \unsure{fidélement|fidèlement} par l'\unsure{ens} $\widetilde F$ des strates de $F$.

\deleted{\unsure{En fait}} \deleted{\illegible{1}} En fait, $\mathcal{M}$ sera muni d'une relation d'ordre $\leq$, cf plus bas, et $\widetilde F\inserted{\subset\mathcal{M}}$ sera une partie fermée de \deleted{\unsure{la}} $\mathcal{M}$, et pour tout $X\in\widetilde F$, \deleted{\unsure{$X$}} on aura
\[
\widetilde X=\{Y\in\mathcal{M}\mid Y\leq\unsure{x|X}\}\tag*{(1.5\unsure{.})}
\]
